\section{Spring Fever Intoxication}

\begin{quotex}
That Self, which is free from sin, free from old age, from death and grief, from hunger and thirst, which desires nothing but what it ought to desire, and imagines nothing but what it ought to imagine, it is that which we should search out; that which we must try to understand. He who reaches that Self and understands it gains all the world of desires.
\flright{\textsc{Chandogya Upanishad}}

For a just man shall fall seven times and shall rise again: but the wicked shall fall down into evil.
\flright{\textsc{Proverbs 24:16}}
\end{quotex}


\paragraph{Random Thoughts}


It is easy to sound intelligent if you make up facts as you go.

People see the tip of the iceberg and think they are brilliant for having noticed it. Even the Captain of the Titanic noticed the tip of the iceberg.

Spiritually empty people often gravitate to partisan politics as an alternative. Once ensconced inside an ideology, they can then identify the orthodox and the heretics, sinners and saints, the saved and the damned. This is an artificial manifestation of \textbf{Carl Schmitt}'s friend-enemy distinction. Meanwhile, the real enemy goes unnoticed.

Recently someone wrote to me that reality is given, and we need to passively know its ``truths''. That can only be the case in a depersonalized cosmos. However, the cosmos is buzzing with persons who are ``willing'' reality. How do we know that? By willing, of course, and the self-understanding of oneself as an ``I'' in the world, not as an ``it''.

\paragraph{Intellectuals and Charisma}


A British poll selected the top 10 public intellectuals\footnote{\url{http://rt.com/uk/244225-brand-influential-thinker-voted}}. Other than Klugman and Habermas, I am not familiar with any of them. My list would be different, assuming I could even reach 10. So I watched a one man show by \textbf{Russell Brand} (\#4) titled \emph{Messiah Complex}, or rather 30 minutes of it. (I assumed it would not get better.) He is just intelligent enough to be able to digest and regurgitate propaganda effectively, but not ``solid food''. His theme is that we project qualities onto certain people, as ``Messiahs''. Interestingly enough he referred to ``Stories we are told'', as the reason for this complex. However, instead of transcending all points of view, he wants to replace those stories with a rather bland left wing story.

As \textbf{Julius Evola} said, people at the first stage simply accept a spontaneously arising world view---or ``story''---as true without question. At the next stage, the first stage is called into question. Hence, the spontaneously arising worldview is replaced with an artificially constructed scientific, ideological, political, or religious worldview. Brand is at a rudimentary second stage whose success is convincing those at the first stage to uncritically adopt his ``story'' based on his fame and charisma.

I was conversing with three Indian colleagues last week, two exoteric and one who is searching. The latter one was quite taken by \textbf{Osho} to my chagrin. I told them the story of Osho here and also in India. I also gave them a little lesson on the Vedanta; since they appreciated it, I asked them to buy me a Maserati, since Osho's followers bought him many cars. They politely declined.

The Osho guy kept insisting on the value in his lectures. I told him I would take an hour and teach him everything Osho knew, provided he would buy me a car. He was interested in the talk, but not in the purchase.

I asked them what the problem is and I got an answer I did not want to hear. They all agreed that, despite the value in my talk, no one would buy me a car because I lacked ``charisma''.

I pondered that until a recent incident on facebook made that crystal clear to me. Someone had posted a photo with Guenon's quote\footnote{\url{https://gornahoor.net/?p=5521}} about the end of the world being the end of an illusion. That got several dozen ``likes'', whereas my many posts on the very same topic might get 4 or 5 ``likes''. That must be due to the lack of charisma. I wonder how many of those likers understand that particular quote; probably no more than one.

\paragraph{Fairy Tales}


I got to watch some ``children's programming'' last weekend. You can see it is TV written by teams of psychologists, to ``tell stories'' forming the first stage. They are artificially created stories, unlike fairy tales, which arise from the soul life of a people.

My psyche was formed by fairy tales and nursery rhymes. How may such stories and poems can you tell by heart? In elementary school, we even learned Greek and Norse myths.

I suspect that those under 35 are formed by cartoons created by psychologists. There were no 30 minute cartoons in my day.

\paragraph{Spiritual Delusion}


In my talk with those Indian colleagues, I told them to distrust any movement that attracts Westerners (like Osho). Americans, especially, believe they are entitled to start at the top. No bhakti or karma yoga for them, but straight to raja yoga. In Buddhism, too, I've noticed, they want to go straight to Dzogchen or Tantra without the years of preparation.

The desire to be considered ``spiritual'' is a tendency, particularly among a certain type of woman, but even some men. They are usually attractive, successful, financially secure, but self-knowledge is not a commodity to be bought and paid for. And a charming personality does not help. Hence they jump at every novelty, speak in elevated tones, but are totally lacking in self-knowledge. They crave a different type of attention. My advice to them is:

\begin{quotex}
Only God my dear, could love you for yourself alone and not your yellow hair.
\flright{\textsc{William Butler Yeats}}
\end{quotex}

Why does Gornahoor repel women? I've met women who have read books, visited ashrams, and even have their own personal shaman or Jungian analyst. I've spoken to Buddhists (which apparently is a polite way to be an atheist these days), artists, doctors, and so on. But as soon as I point them to Gornahoor, they shut down. My sister calls me a ``shaman'', so perhaps I should follow that path. What should I charge?

Recently, I spoke to a woman who self identifies as a ``very conservative Christian''. Curiously, she ranks \emph{Fifty Shades of Grey} as one of her favorite books, for its ``honesty''. Probably it reminds her of Paul's description of a Christian marriage. Always aiming to please, I asked her if she would like to try a relationship like that. She hung up on me. Just a warning to you guys: it may sound good on paper, but you probably won't want a slave underfoot all the time. Unless you can afford to give her a credit card to keep her out of the house shopping, give it a pass.

\paragraph{Life Changing Events}


I subscribe to a few ``self help'' mailing lists to take the pulse on things. I hear about all these ``amazing'' meetings with high-charged people. They are pumped up at the end. I don't know why, but that seems to be worth hundreds, or even thousands, of dollars.

The Gnosis group is not like that, no one is ``amazing''. It is hard work, but a beautiful work.

One cute and perky woman sends me weekly notices about ``life changing'' opportunities. Apparently, there is a market niche based on people who want to change their lives. (She claims to be earning a 7 figure income.)

I personally don't grok the appeal of that. If I kept changing my life that often, by this point I would have been a zoo keeper, a dance instructor on a cruise line, and would probably now be singing show tunes at a transvestite bar. No thanks, I'll keep the life I have.

\paragraph{Sequel}


This was supposed to be a post about self-mastery, but I ran out of space. As I've insisted, a man should live the philosophy he holds. Hence, you won't find, in my case, some wonderful conversion story but rather a series of attempts at self-mastery, of the subtle struggling to rule the dense. That will come soon. But first, we need the promised story of the third migration from Aeneas, but he isn't returning my phone calls.

On the path of gnosis, you will be freed of anxiety, addictions, self-doubts, and lack of confidence. What is that worth to you?


\flrightit{Posted on 2015-03-29 by Cologero}

\begin{center}* * *\end{center}

\begin{footnotesize}
\begin{sffamily}
\texttt{Ronin Yogi on 2015-03-29 at 23:45 said:}

Fine thoughts today. Gnosis is of much ``worth''.

\hfill

\texttt{Wayne Ferguson on 2015-03-30 at 01:31 said:} 

``Hence, you won't find, in my case, some wonderful conversion story but rather a series of attempts at self-mastery, of the subtle struggling to rule the dense.''

Would it be more accurate to say that the dense has been struggling to rule the dense instead of submitting to the subtle? Think also of the conflict in Romans 7 giving way to the life of the Spirit in Romans 8.

Quoting Letter XI of Tomberg's ``Meditations on the Tarot'':

(This thing is the strongest of all powers,

the force of all forces,

for it overcometh every subtle thing

and doth penetrate every solid substance.)

(tabula Smaragdina, 9)

``Virgin Nature therefore has her part in all the miracles. And it is virgin Nature participating actively in the miracles of divine magic which is the subject of the eleventh Arcanum of the Tarot, force, representing a woman victorious over a lion, holding its jaws open with her hands. The woman does so with the same apparent ease---without effort ---with which the Magician of the first Arcanum handles his objects. Moreover, she wears a hat similar to that of the Magician ---in the form of a lemniscate. One could say that the two stand equally under the sign of rhythm ---the respiration of eternity---the sign ? ; and that the two manifest two aspects of a single principle, namely that effort signifies the presence of an obstacle, whilst natural integrity on the one hand, and undivided attention on the other hand, exclude inner conflict---and therefore every obstacle, and therefore all effort. Just as perfect concentration takes place effortlessly, so does true force act without effort. Now, the Magician is the Arcanum of the wholeness of consciousness, or concentration without effort; Force is the Arcanum of the natural integrity of being, or power without effort. Because force subdues the lion not by force similar to that of the lion, but rather by force of a higher order and on a higher plane. This is the Arcanum of Force'' (Letter 11, Force).

\hfill

\texttt{Desert Scout on 2015-03-30 at 11:00 said:}

``My sister calls me a ``shaman'', so perhaps I should follow that path. What should I charge?''

My rate is \$65 an hour, cash, and the girls typically tip. When the government is footing the bill charge \$240.

Before I found Gornahoor I was suicidalily depressed. Thanks to you I discovered the writings of Tomberg, and through converting to and living the Sacramental life of the Catholic faith I've been freed from my anxiety, depression, and suicidal ideation. One of the problems in any field, both self help and medicine, is that if you actually cure the client you put yourself out of business. It's not uncommon that I have a patient cancel their follow up appointment because they're out of pain.

For what it's worth, you've made a difference in at least one person's life. When you go before the Almighty you can at least say you saved one of the Master's wayward sheep. I don't know if that's what you aimed for, but it's something that's hopefully worthwhile.

\hfill

\texttt{rhondda9 on 2015-03-30 at 12:06 said:}

It is very hard for women to admit that everything they think they know is a lie and that they have to start at the beginning again. It is really a terrifying realization and better to pretend you never thought it.

\hfill

\texttt{Avery on 2015-03-31 at 16:50 said:}

As recently as two years ago, people were still writing blogs with little to no evidence that anyone was reading at all. Now, it is impossible to avoid the quantification of the power your words place on others --- what is called without irony ``metrics''. Usually, the metric is the number of advertisements viewed, but it is impossible to avoid the metric of Facebook likes even for Gornahoor.

I can't remember whether I clicked like for either of your posts. Of course, the reality of Guenon's message is the same whether or not people like, or ``like'' it. There is really an unbelievable destruction of thought going on with this infiltration of metrics everywhere.

\hfill

\texttt{Paulo on 2015-03-31 at 17:57 said:}

Fear your captors,fear your masters-Juan Matus

\hfill

\texttt{Cologero on 2015-03-31 at 19:06 said:}

You need to develop a sensayuma, Avery, although that is not quantifiable.

\hfill

\texttt{Jacob on 2015-04-01 at 12:05 said:}

Yea Avery. I guess we could say another sign of the reign of quantity?

\hfill

\texttt{Paulo on 2015-04-01 at 12:19 said:}

``It is very hard for women to admit that everything they think they know is a lie and that they have to start at the beginning again. It is really a terrifying realization and better to pretend you never thought it.''

If there is a real man there in the moment,he will cover that situation too!

\hfill

\texttt{rhondda9 on 2015-04-01 at 13:23 said:}

Was that your attempt at humour Avery?

\hfill

\texttt{rhondda9 on 2015-04-01 at 15:11 said:}

Oops Paulo, I mean. Back to the nunnery for me.

\hfill

\texttt{Paul on 2015-06-12 at 23:59 said:}

``Americans, especially, believe they are entitled to start at the top. No bhakti or karma yoga for them, but straight to raja yoga. In Buddhism, too, I've noticed, they want to go straight to Dzogchen or Tantra without the years of preparation... ''

Agreed in general, although some additional observations may be germane, especially with respect to Dzogchen. As a student of Tibetan scholar Namkhai Norbu, I have come to realize that given the situation in Tibet throughout the last century or so, there would be very little chance of any Tibetans practicing Dzogchen or indeed any kind of tantra openly under Chinese eyes, without persecution. Also, Dzogchen itself has not been taught to the average Tibetan Vajrayana practitioner, as such an ``advanced'' type of practice is well beyond the means of those who are only at the level of mantra recitations and prostrations; therefore, it has always been more or less an underground path, transmitted secretly to those who are ``ready.''

Norbu has also alluded that those who approach this path here in the West already have past-life connections to it since, given the above, most practitioners who were ready to approach or continue this clandestine path of development (as Tibetans in previous lives) would of necessity need to be reborn in the West to obtain the freedom to avail themselves of that opportunity. This makes perfect sense.

\end{sffamily}
\end{footnotesize}

